\documentclass[
  man,
  floatsintext,
  longtable,
  nolmodern,
  notxfonts,
  notimes,
  draftfirst,
  colorlinks=true,linkcolor=blue,citecolor=blue,urlcolor=blue]{apa7}

\usepackage{amsmath}
\usepackage{amssymb}




\RequirePackage{longtable}
\RequirePackage{threeparttablex}

\makeatletter
\renewcommand{\paragraph}{\@startsection{paragraph}{4}{\parindent}%
	{0\baselineskip \@plus 0.2ex \@minus 0.2ex}%
	{-.5em}%
	{\normalfont\normalsize\bfseries\typesectitle}}

\renewcommand{\subparagraph}[1]{\@startsection{subparagraph}{5}{0.5em}%
	{0\baselineskip \@plus 0.2ex \@minus 0.2ex}%
	{-\z@\relax}%
	{\normalfont\normalsize\bfseries\itshape\hspace{\parindent}{#1}\textit{\addperi}}{\relax}}
\makeatother




\usepackage{longtable, booktabs, multirow, multicol, colortbl, hhline, caption, array, float, xpatch}
\setcounter{topnumber}{2}
\setcounter{bottomnumber}{2}
\setcounter{totalnumber}{4}
\renewcommand{\topfraction}{0.85}
\renewcommand{\bottomfraction}{0.85}
\renewcommand{\textfraction}{0.15}
\renewcommand{\floatpagefraction}{0.7}

\usepackage{tcolorbox}
\tcbuselibrary{listings,theorems, breakable, skins}
\usepackage{fontawesome5}

\definecolor{quarto-callout-color}{HTML}{909090}
\definecolor{quarto-callout-note-color}{HTML}{0758E5}
\definecolor{quarto-callout-important-color}{HTML}{CC1914}
\definecolor{quarto-callout-warning-color}{HTML}{EB9113}
\definecolor{quarto-callout-tip-color}{HTML}{00A047}
\definecolor{quarto-callout-caution-color}{HTML}{FC5300}
\definecolor{quarto-callout-color-frame}{HTML}{ACACAC}
\definecolor{quarto-callout-note-color-frame}{HTML}{4582EC}
\definecolor{quarto-callout-important-color-frame}{HTML}{D9534F}
\definecolor{quarto-callout-warning-color-frame}{HTML}{F0AD4E}
\definecolor{quarto-callout-tip-color-frame}{HTML}{02B875}
\definecolor{quarto-callout-caution-color-frame}{HTML}{FD7E14}

%\newlength\Oldarrayrulewidth
%\newlength\Oldtabcolsep


\usepackage{hyperref}




\providecommand{\tightlist}{%
  \setlength{\itemsep}{0pt}\setlength{\parskip}{0pt}}
\usepackage{longtable,booktabs,array}
\usepackage{calc} % for calculating minipage widths
% Correct order of tables after \paragraph or \subparagraph
\usepackage{etoolbox}
\makeatletter
\patchcmd\longtable{\par}{\if@noskipsec\mbox{}\fi\par}{}{}
\makeatother
% Allow footnotes in longtable head/foot
\IfFileExists{footnotehyper.sty}{\usepackage{footnotehyper}}{\usepackage{footnote}}
\makesavenoteenv{longtable}

\usepackage{graphicx}
\makeatletter
\newsavebox\pandoc@box
\newcommand*\pandocbounded[1]{% scales image to fit in text height/width
  \sbox\pandoc@box{#1}%
  \Gscale@div\@tempa{\textheight}{\dimexpr\ht\pandoc@box+\dp\pandoc@box\relax}%
  \Gscale@div\@tempb{\linewidth}{\wd\pandoc@box}%
  \ifdim\@tempb\p@<\@tempa\p@\let\@tempa\@tempb\fi% select the smaller of both
  \ifdim\@tempa\p@<\p@\scalebox{\@tempa}{\usebox\pandoc@box}%
  \else\usebox{\pandoc@box}%
  \fi%
}
% Set default figure placement to htbp
\def\fps@figure{htbp}
\makeatother


% definitions for citeproc citations
\NewDocumentCommand\citeproctext{}{}
\NewDocumentCommand\citeproc{mm}{%
  \begingroup\def\citeproctext{#2}\cite{#1}\endgroup}
\makeatletter
 % allow citations to break across lines
 \let\@cite@ofmt\@firstofone
 % avoid brackets around text for \cite:
 \def\@biblabel#1{}
 \def\@cite#1#2{{#1\if@tempswa , #2\fi}}
\makeatother
\newlength{\cslhangindent}
\setlength{\cslhangindent}{1.5em}
\newlength{\csllabelwidth}
\setlength{\csllabelwidth}{3em}
\newenvironment{CSLReferences}[2] % #1 hanging-indent, #2 entry-spacing
 {\begin{list}{}{%
  \setlength{\itemindent}{0pt}
  \setlength{\leftmargin}{0pt}
  \setlength{\parsep}{0pt}
  % turn on hanging indent if param 1 is 1
  \ifodd #1
   \setlength{\leftmargin}{\cslhangindent}
   \setlength{\itemindent}{-1\cslhangindent}
  \fi
  % set entry spacing
  \setlength{\itemsep}{#2\baselineskip}}}
 {\end{list}}
\usepackage{calc}
\newcommand{\CSLBlock}[1]{\hfill\break\parbox[t]{\linewidth}{\strut\ignorespaces#1\strut}}
\newcommand{\CSLLeftMargin}[1]{\parbox[t]{\csllabelwidth}{\strut#1\strut}}
\newcommand{\CSLRightInline}[1]{\parbox[t]{\linewidth - \csllabelwidth}{\strut#1\strut}}
\newcommand{\CSLIndent}[1]{\hspace{\cslhangindent}#1}





\usepackage{newtx}

\defaultfontfeatures{Scale=MatchLowercase}
\defaultfontfeatures[\rmfamily]{Ligatures=TeX,Scale=1}





\title{Intrusive and Voluntary Memory from a Single System: A
Retrieved-Context Account of Selective Interference}


\shorttitle{Intrusive and Voluntary Memory from a Single System}


\usepackage{etoolbox}






\author{Jordan B Gunn}



\affiliation{
{Department of Psychology, University of Cambridge}}




\leftheader{Gunn}



\abstract{Visuospatial tasks performed after trauma-film encoding, or
after a later reminder of the film, can reduce intrusive memories while
leaving voluntary memory relatively intact. This selective interference
effect is often treated as evidence that intrusive and voluntary
remembering depend on differently vulnerable memory representations.
Here we simulate an alternative retrieved-context account using an
emotional Context Maintenance and Retrieval (eCMR) model, asking whether
context reinstatement, competitor learning, and retrieval control can
generate the selective-interference pattern without separately
vulnerable representations. In the model, film context remains active
after encoding and can later be reinstated by reminders, so an
interfering task performed in either state binds task representations to
overlapping context. When later retrieval contexts cue both film
representations and task competitors, uncontrolled retrieval from
context cues is less likely to sample trauma-film items in
intrusion-like retrieval, even though film-item strength is not reduced.
Deliberate recall is less vulnerable when task goals and control
processes reinstate film context and prioritize film items over
competitors for retrieval; film-specific cues can also steer retrieval
toward film context across test formats. Together, the simulations
reproduce selective impairment of intrusion-like retrieval and specify
when the intrusion-voluntary-memory dissociation should be more or less
pronounced, helping to interpret heterogeneity across empirical
paradigms. Selective interference is therefore compatible with a single
retrieved-context system and is not by itself diagnostic of separate
representational systems or selective trace weakening. }

\keywords{intrusive memories, voluntary memory, trauma-film
paradigm, visuospatial intervention, selective
interference, retrieved-context theory, computational modeling}

\authornote{ 

\par{  This is an in-progress working draft. Please do not cite, quote,
or circulate without the authors' written permission.     }
}

\makeatletter
\let\endoldlt\endlongtable
\def\endlongtable{
\hline
\endoldlt
}
\makeatother

\urlstyle{same}



\makeatletter
\@ifpackageloaded{caption}{}{\usepackage{caption}}
\AtBeginDocument{%
\ifdefined\contentsname
  \renewcommand*\contentsname{Table of contents}
\else
  \newcommand\contentsname{Table of contents}
\fi
\ifdefined\listfigurename
  \renewcommand*\listfigurename{List of Figures}
\else
  \newcommand\listfigurename{List of Figures}
\fi
\ifdefined\listtablename
  \renewcommand*\listtablename{List of Tables}
\else
  \newcommand\listtablename{List of Tables}
\fi
\ifdefined\figurename
  \renewcommand*\figurename{Figure}
\else
  \newcommand\figurename{Figure}
\fi
\ifdefined\tablename
  \renewcommand*\tablename{Table}
\else
  \newcommand\tablename{Table}
\fi
}
\@ifpackageloaded{float}{}{\usepackage{float}}
\floatstyle{ruled}
\@ifundefined{c@chapter}{\newfloat{codelisting}{h}{lop}}{\newfloat{codelisting}{h}{lop}[chapter]}
\floatname{codelisting}{Listing}
\newcommand*\listoflistings{\listof{codelisting}{List of Listings}}
\makeatother
\makeatletter
\makeatother
\makeatletter
\@ifpackageloaded{caption}{}{\usepackage{caption}}
\@ifpackageloaded{subcaption}{}{\usepackage{subcaption}}
\makeatother

% From https://tex.stackexchange.com/a/645996/211326
%%% apa7 doesn't want to add appendix section titles in the toc
%%% let's make it do it
\makeatletter
\xpatchcmd{\appendix}
  {\par}
  {\addcontentsline{toc}{section}{\@currentlabelname}\par}
  {}{}
\makeatother

%% Disable longtable counter
%% https://tex.stackexchange.com/a/248395/211326

\usepackage{etoolbox}

\makeatletter
\patchcmd{\LT@caption}
  {\bgroup}
  {\bgroup\global\LTpatch@captiontrue}
  {}{}
\patchcmd{\longtable}
  {\par}
  {\par\global\LTpatch@captionfalse}
  {}{}
\apptocmd{\endlongtable}
  {\ifLTpatch@caption\else\addtocounter{table}{-1}\fi}
  {}{}
\newif\ifLTpatch@caption
\makeatother

\begin{document}

\maketitle


\setcounter{secnumdepth}{-\maxdimen} % remove section numbering

\setlength\LTleft{0pt}


\section{Introduction}\label{introduction}

Traumatic experience foregrounds a tension between durability and
selectivity in episodic memory. Highly consequential events should be
retained well enough in memory that critical details remain accessible
long into the future. At the same time, adaptive retrieval is selective:
what comes to mind should normally depend on current cues, goals, and
safety. Intrusive memories after trauma suggest that durable encoding
can outpace regulatory control. The event returns without being
deliberately sought and without being well matched to the demands of the
present (\citeproc{ref-brewin2010intrusive}{Brewin et al., 2010};
\citeproc{ref-ehlers2000cognitive}{Ehlers \& Clark, 2000};
\citeproc{ref-iyadurai2019intrusive}{Iyadurai et al., 2019};
\citeproc{ref-krans2009intrusive}{Krans et al., 2009}). Understanding
why such memories arise, and how they differ from voluntary memory for
the same event, is therefore a shared target for both clinical science
and basic memory theory.

The trauma-film paradigm provides a controlled setting for comparing
intrusive and voluntary memory for the same event
(\citeproc{ref-holmes2008inducing}{Holmes \& Bourne, 2008};
\citeproc{ref-james2016trauma}{James et al., 2016}). Participants view
distressing film material, record later intrusive memories of the film,
and often complete voluntary-memory tests for the same material. This
design makes it possible to ask whether the same encoded event can
remain available to deliberate retrieval while becoming more or less
likely to return involuntarily. In such research, voluntary memory can
be relatively preserved, and its relationship to intrusion frequency can
be weak, such that the two measures dissociate
(\citeproc{ref-holmes2004trauma}{Holmes et al., 2004};
\citeproc{ref-james2016trauma}{James et al., 2016}). The same encoded
episode can therefore support both uncontrolled and controlled access,
and these modes can come apart empirically
(\citeproc{ref-holmes2004trauma}{Holmes et al., 2004};
\citeproc{ref-iyadurai2019intrusive}{Iyadurai et al., 2019}).

Visuospatial interference tasks make this dissociation experimentally
sharper. When brief tasks such as Tetris are performed after film
encoding or after a later film reminder, subsequent intrusions can be
reduced while voluntary memory for the same material shows little
measurable impairment
(\citeproc{ref-asselbergs2023systematic}{Asselbergs et al., 2023};
\citeproc{ref-deeprose2012imagery}{Deeprose et al., 2012};
\citeproc{ref-holmes2009can}{Holmes et al., 2009},
\citeproc{ref-holmes2010key}{2010};
\citeproc{ref-james2015computer}{James et al., 2015};
\citeproc{ref-kessler2020visuospatial}{Kessler et al., 2020};
\citeproc{ref-lau2019intrusive}{Lau-Zhu et al., 2019}). Even for
memories of real-life traumatic events, including motor vehicle
accidents, emergency caesarean section, and war-related trauma, the same
post-reminder intervention can reduce intrusions while leaving
deliberate recall relatively intact
(\citeproc{ref-horsch2017reducing}{Horsch et al., 2017};
\citeproc{ref-iyadurai2018preventing}{Iyadurai et al., 2018};
\citeproc{ref-kanstrup2021single}{Kanstrup et al., 2021}). At the same
time, the evidence base is heterogeneous, and recent reviews and
replications caution against treating the effect as settled or uniform
(\citeproc{ref-asselbergs2023systematic}{Asselbergs et al., 2023};
\citeproc{ref-varma2024experimental}{Varma et al., 2024};
\citeproc{ref-wessel2025evidence}{Wessel et al., 2025}). This selective
interference effect therefore raises a theoretical question: when and
why can a post-encoding or post-reminder manipulation alter one mode of
retrieval more than another, even while both draw on the same recent
experience?

One prominent interpretation treats the selective interference effect as
evidence for separable memory representations
(\citeproc{ref-brewin1996dual}{Brewin et al., 1996};
\citeproc{ref-brewin2014episodic}{Brewin, 2014}). Under this
dual-representation account, traumatic experience produces distinct
sensory and contextual memory traces --- the first supporting
involuntary re-experiencing, the second supporting deliberate recall
(\citeproc{ref-brewin2010intrusive}{Brewin et al., 2010};
\citeproc{ref-brewin2014episodic}{Brewin, 2014}). A visuospatial
interference task competes for the perceptual processing resources
needed to consolidate the sensory trace, reducing later intrusions while
leaving the contextual trace and the voluntary memory it supports
relatively intact (\citeproc{ref-holmes2009can}{Holmes et al., 2009};
\citeproc{ref-james2015computer}{James et al., 2015}). When the
intervention follows a reminder at longer delays, the account invokes
reconsolidation: the reminder reactivates the sensory trace, and the
interfering task disrupts its restabilization during a time-limited
window (\citeproc{ref-james2015computer}{James et al., 2015};
\citeproc{ref-kindt2009beyond}{Kindt et al., 2009}). The account thus
ties selective interference to separate memory representations,
modality-specific disruption of one trace, and a discrete
reconsolidation window.

Here we argue that selective interference need not be interpreted as
evidence for differently vulnerable memory representations. It may
instead reflect differences in retrieval processes operating over a
shared set of episodic representations. Mechanisms supporting selective,
goal-directed retrieval in voluntary recall may be sufficient to
sidestep or overcome interference that nonetheless reduces the
likelihood of unintentional recall. On this interpretation, the
selective-interference effect reflects a capacity to intentionally focus
retrieval on goal-relevant event information while reducing the
influence of competing memories.

We formalize and evaluate this proposal within a retrieved-context
framework drawing on the emotional Context Maintenance and Retrieval
model family, including eCMR and CMR3
(\citeproc{ref-cohen2022memory}{Cohen \& Kahana, 2022};
\citeproc{ref-talmi2019retrieved}{Talmi et al., 2019}).
Retrieved-context models explain recall by treating context as a
gradually changing state that is updated by studied items, reinstated by
recalled items, and used to cue later memory search
(\citeproc{ref-howard2002distributed}{Howard \& Kahana, 2002};
\citeproc{ref-polyn2009context}{Polyn et al., 2009}). Emotional
extensions adds source-like emotional context and emotion-dependent
binding strength, allowing the same architecture to capture how
emotional content changes the organization and success of recall Cohen
\& Kahana (\citeproc{ref-cohen2022memory}{2022}). Variants also model
strategic control over memory search, allowing the model to capture how
retrieval goals and task demands to shape recall
(\citeproc{ref-cohen2022memory}{Cohen \& Kahana, 2022};
\citeproc{ref-lohnas2015expanding}{Lohnas et al., 2015};
\citeproc{ref-sederberg2008context}{Sederberg et al., 2008}). Together,
these commitments make the framework a principled setting for testing
whether selective interference can arise within a single emotional
episodic-memory system.

Consolidation and reconsolidation accounts locate selective interference
in a change to the availability of the film representation after
encoding or reminder reactivation. A context-binding account shifts the
explanatory focus to the retrieval consequences of encoding later
experiences in contexts that overlap with the original event
(\citeproc{ref-yonelinas2019contextual}{Yonelinas et al., 2019}). For
selective interference, our model treats the original event, any
reminder, and subsequent interference as episodes encoded by binding
item representations to a continuously evolving context representation
(\citeproc{ref-howard2002distributed}{Howard \& Kahana, 2002};
\citeproc{ref-polyn2009context}{Polyn et al., 2009};
\citeproc{ref-yonelinas2019contextual}{Yonelinas et al., 2019}). Later
retrieval is driven by the current context state as a cue, so
competition depends on how strongly different episodes are associated
with overlapping regions of context. Deliberate recall is more resistant
because it combines a more target-oriented initial cue with a sharper
decision rule, both of which favor film items over competitors. The
dissociation between involuntary and voluntary recall is thus attributed
to differences in retrieval state applied to the same stored episodes.

\section{Empirical and Theoretical
Target}\label{empirical-and-theoretical-target}

\section{A Retrieved-Context Account of Selective
Interference}\label{a-retrieved-context-account-of-selective-interference}

\section{Model and Paradigm}\label{model-and-paradigm}

\section{Simulation Overview}\label{simulation-overview}

\section{Results}\label{results}

\subsection{Simulation 1: Selective Interference Across Retrieval
Configurations}\label{simulation-1-selective-interference-across-retrieval-configurations}

\subsection{Simulation 2: Competitor Encoding in Reinstated
Context}\label{simulation-2-competitor-encoding-in-reinstated-context}

\subsection{Simulation 3: Contextual Overlap Versus Generic Task
Load}\label{simulation-3-contextual-overlap-versus-generic-task-load}

\subsection{Simulation 4: Test-Format
Heterogeneity}\label{simulation-4-test-format-heterogeneity}

\subsection{Robustness and Diagnostic
Checks}\label{robustness-and-diagnostic-checks}

\section{General Discussion}\label{general-discussion}

\section{References}\label{references}

\phantomsection\label{refs}
\begin{CSLReferences}{1}{0}
\bibitem[\citeproctext]{ref-asselbergs2023systematic}
Asselbergs, J., van Bentum, J., Riper, H., Cuijpers, P., Holmes, E. A.,
\& Sijbrandij, M. (2023). A systematic review and meta-analysis of the
effect of cognitive interventions to prevent intrusive memories using
the trauma film paradigm. \emph{Journal of Psychiatric Research},
\emph{159}, 116--129.
\url{https://doi.org/10.1016/j.jpsychires.2023.01.028}

\bibitem[\citeproctext]{ref-brewin2014episodic}
Brewin, C. R. (2014). Episodic memory, perceptual memory, and their
interaction: Foundations for a theory of posttraumatic stress disorder.
\emph{Psychological Bulletin}, \emph{140}(1), 69--97.
\url{https://doi.org/10.1037/a0033722}

\bibitem[\citeproctext]{ref-brewin1996dual}
Brewin, C. R., Dalgleish, T., \& Joseph, S. (1996). A dual
representation theory of posttraumatic stress disorder.
\emph{Psychological Review}, \emph{103}(4), 670--686.
\url{https://doi.org/10.1037/0033-295X.103.4.670}

\bibitem[\citeproctext]{ref-brewin2010intrusive}
Brewin, C. R., Gregory, J. D., Lipton, M., \& Burgess, N. (2010).
Intrusive images in psychological disorders: Characteristics, neural
mechanisms, and treatment implications. \emph{Psychological Review},
\emph{117}(1), 210--232. \url{https://doi.org/10.1037/a0018113}

\bibitem[\citeproctext]{ref-cohen2022memory}
Cohen, R. T., \& Kahana, M. J. (2022). A memory-based theory of
emotional disorders. \emph{Psychological Review}, \emph{129}(4),
742--776. \url{https://doi.org/10.1037/rev0000334}

\bibitem[\citeproctext]{ref-deeprose2012imagery}
Deeprose, C., Zhang, S., DeJong, H., Dalgleish, T., \& Holmes, E. A.
(2012). Imagery in the aftermath of viewing a traumatic film: Using
cognitive tasks to modulate the development of involuntary memory.
\emph{Journal of Behavior Therapy and Experimental Psychiatry},
\emph{43}(2), 758--764.
\url{https://doi.org/10.1016/j.jbtep.2011.10.008}

\bibitem[\citeproctext]{ref-ehlers2000cognitive}
Ehlers, A., \& Clark, D. M. (2000). A cognitive model of posttraumatic
stress disorder. \emph{Behaviour Research and Therapy}, \emph{38}(4),
319--345. \url{https://doi.org/10.1016/S0005-7967(99)00123-0}

\bibitem[\citeproctext]{ref-holmes2008inducing}
Holmes, E. A., \& Bourne, C. (2008). Inducing and modulating intrusive
emotional memories: A review of the trauma film paradigm. \emph{Acta
Psychologica}, \emph{127}(3), 553--566.
\url{https://doi.org/10.1016/j.actpsy.2007.11.002}

\bibitem[\citeproctext]{ref-holmes2004trauma}
Holmes, E. A., Brewin, C. R., \& Hennessy, R. G. (2004). Trauma films,
information processing, and intrusive memory development. \emph{Journal
of Experimental Psychology: General}, \emph{133}(1), 3--22.
\url{https://doi.org/10.1037/0096-3445.133.1.3}

\bibitem[\citeproctext]{ref-holmes2009can}
Holmes, E. A., James, E. L., Coode-Bate, T., \& Deeprose, C. (2009). Can
playing the computer game {Tetris} reduce the build-up of flashbacks for
trauma? A proposal from cognitive science. \emph{PLoS ONE}, \emph{4}(1),
e4153. \url{https://doi.org/10.1371/journal.pone.0004153}

\bibitem[\citeproctext]{ref-holmes2010key}
Holmes, E. A., James, E. L., Kilford, E. J., \& Deeprose, C. (2010). Key
steps in developing a cognitive vaccine against traumatic flashbacks:
Visuospatial {Tetris} versus verbal {Pub Quiz}. \emph{PLoS ONE},
\emph{5}(11), e13706. \url{https://doi.org/10.1371/journal.pone.0013706}

\bibitem[\citeproctext]{ref-horsch2017reducing}
Horsch, A., Vial, Y., Favrod, C., Harari, M. M., Blackwell, S. E.,
Watson, P., Iyadurai, L., Bonsall, M. B., \& Holmes, E. A. (2017).
Reducing intrusive traumatic memories after emergency caesarean section:
A proof-of-principle randomized controlled study. \emph{Behaviour
Research and Therapy}, \emph{94}, 36--47.
\url{https://doi.org/10.1016/j.brat.2017.03.018}

\bibitem[\citeproctext]{ref-howard2002distributed}
Howard, M. W., \& Kahana, M. J. (2002). A distributed representation of
temporal context. \emph{Journal of Mathematical Psychology},
\emph{46}(3), 269--299. \url{https://doi.org/10.1006/jmps.2001.1388}

\bibitem[\citeproctext]{ref-iyadurai2018preventing}
Iyadurai, L., Blackwell, S. E., Meiser-Stedman, R., Watson, P. C.,
Bonsall, M. B., Geddes, J. R., Nobre, A. C., \& Holmes, E. A. (2018).
Preventing intrusive memories after trauma via a brief intervention
involving {Tetris} computer game play in the emergency department: A
proof-of-concept randomized controlled trial. \emph{Molecular
Psychiatry}, \emph{23}(3), 674--682.
\url{https://doi.org/10.1038/mp.2017.23}

\bibitem[\citeproctext]{ref-iyadurai2019intrusive}
Iyadurai, L., Visser, R. M., Lau-Zhu, A., Porcheret, K., Horsch, A.,
Holmes, E. A., \& James, E. L. (2019). Intrusive memories of trauma: A
target for research bridging cognitive science and its clinical
application. \emph{Clinical Psychology Review}, \emph{69}, 67--82.
\url{https://doi.org/10.1016/j.cpr.2018.08.005}

\bibitem[\citeproctext]{ref-james2015computer}
James, E. L., Bonsall, M. B., Hoppitt, L., Tunbridge, E. M., Geddes, J.
R., Milton, A. L., \& Holmes, E. A. (2015). Computer game play reduces
intrusive memories of experimental trauma via reconsolidation-update
mechanisms. \emph{Psychological Science}, \emph{26}(8), 1201--1215.
\url{https://doi.org/10.1177/0956797615583071}

\bibitem[\citeproctext]{ref-james2016trauma}
James, E. L., Lau-Zhu, A., Clark, I. A., Visser, R. M., Hagenaars, M.
A., \& Holmes, E. A. (2016). The trauma film paradigm as an experimental
psychopathology model of psychological trauma: Intrusive memories and
beyond. \emph{Clinical Psychology Review}, \emph{47}, 106--142.
\url{https://doi.org/10.1016/j.cpr.2016.04.010}

\bibitem[\citeproctext]{ref-kanstrup2021single}
Kanstrup, M., Kontio, E., Geranmayeh, A., Olofsdotter Lauri, K., Moulds,
M. L., \& Holmes, E. A. (2021). A single case series using visuospatial
task interference to reduce the number of visual intrusive memories of
trauma with refugees. \emph{Clinical Psychology \& Psychotherapy},
\emph{28}(1), 109--123. \url{https://doi.org/10.1002/cpp.2489}

\bibitem[\citeproctext]{ref-kessler2020visuospatial}
Kessler, H., Schmidt, A.-C., James, E. L., Blackwell, S. E., von
Rauchhaupt, M., Harren, K., Kehyayan, A., Clark, I. A., Sauvage, M.,
Herpertz, S., Axmacher, N., \& Holmes, E. A. (2020). Visuospatial
computer game play after memory reminder delivered three days after a
traumatic film reduces the number of intrusive memories of the
experimental trauma. \emph{Journal of Behavior Therapy and Experimental
Psychiatry}, \emph{67}, 101454.
\url{https://doi.org/10.1016/j.jbtep.2019.01.006}

\bibitem[\citeproctext]{ref-kindt2009beyond}
Kindt, M., Soeter, M., \& Vervliet, B. (2009). Beyond extinction:
Erasing human fear responses and preventing the return of fear.
\emph{Nature Neuroscience}, \emph{12}(3), 256--258.
\url{https://doi.org/10.1038/nn.2271}

\bibitem[\citeproctext]{ref-krans2009intrusive}
Krans, J., Näring, G., Becker, E. S., \& Holmes, E. A. (2009). Intrusive
trauma memory: A review and functional analysis. \emph{Applied Cognitive
Psychology}, \emph{23}(8), 1076--1088.
\url{https://doi.org/10.1002/acp.1611}

\bibitem[\citeproctext]{ref-lau2019intrusive}
Lau-Zhu, A., Henson, R. N., \& Holmes, E. A. (2019). Intrusive memories
and voluntary memory of a trauma film: Differential effects of a
cognitive interference task after encoding. \emph{Journal of
Experimental Psychology: General}, \emph{148}(12), 2154--2180.
\url{https://doi.org/10.1037/xge0000598}

\bibitem[\citeproctext]{ref-lohnas2015expanding}
Lohnas, L. J., Polyn, S. M., \& Kahana, M. J. (2015). Expanding the
scope of memory search: Modeling intralist and interlist effects in free
recall. \emph{Psychological Review}, \emph{122}(2), 337--363.
\url{https://doi.org/10.1037/a0039036}

\bibitem[\citeproctext]{ref-polyn2009context}
Polyn, S. M., Norman, K. A., \& Kahana, M. J. (2009). A context
maintenance and retrieval model of organizational processes in free
recall. \emph{Psychological Review}, \emph{116}(1), 129--156.
\url{https://doi.org/10.1037/a0014420}

\bibitem[\citeproctext]{ref-sederberg2008context}
Sederberg, P. B., Howard, M. W., \& Kahana, M. J. (2008). A
context-based theory of recency and contiguity in free recall.
\emph{Psychological Review}, \emph{115}(4), 893--912.
\url{https://doi.org/10.1037/a0013396}

\bibitem[\citeproctext]{ref-talmi2019retrieved}
Talmi, D., Lohnas, L. J., \& Daw, N. D. (2019). A retrieved context
model of the emotional modulation of memory. \emph{Psychological
Review}, \emph{126}(4), 455--485.
\url{https://doi.org/10.1037/rev0000132}

\bibitem[\citeproctext]{ref-varma2024experimental}
Varma, M. M., Zeng, S., Singh, L., Holmes, E. A., Huang, J., Chiu, M.
H., \& Hu, X. (2024). A systematic review and meta-analysis of
experimental methods for modulating intrusive memories following
lab-analogue trauma exposure in non-clinical populations. \emph{Nature
Human Behaviour}, \emph{8}(10), 1968--1987.
\url{https://doi.org/10.1038/s41562-024-01956-y}

\bibitem[\citeproctext]{ref-wessel2025evidence}
Wessel, I., Krans, J., Albers, C. J., Chauhan, N., Ferreira de Sá, D.
S., Hauck, A., Jaswetz, L., Michael, T., Nixon, R. D. V., Pearson, D.
G., Slattery, I., Smeets, T., Takarangi, M. K. T., Willems, T. R. A., \&
van Schie, K. (2025). Evidence that {Tetris} reduces immediate but not
subsequent daily intrusions of a trauma film: A multilab replication
study. \emph{Collabra: Psychology}, \emph{11}(1), 130791.
\url{https://doi.org/10.1525/collabra.130791}

\bibitem[\citeproctext]{ref-yonelinas2019contextual}
Yonelinas, A. P., Ranganath, C., Ekstrom, A. D., \& Wiltgen, B. J.
(2019). A contextual binding theory of episodic memory: Systems
consolidation reconsidered. \emph{Nature Reviews Neuroscience},
\emph{20}(6), 364--375. \url{https://doi.org/10.1038/s41583-019-0150-4}

\end{CSLReferences}






\end{document}
